\documentclass[11pt]{article}
\usepackage[margin=1in]{geometry}
\usepackage{amsmath,amssymb,amsthm,booktabs,enumitem,hyperref}
\hypersetup{colorlinks=true,linkcolor=blue,urlcolor=blue,citecolor=blue}
\newtheorem{theorem}{Theorem}
\newcommand{\RLM}{\operatorname{RLMv}}
\newcommand{\EXC}{\operatorname{EXCv}}
\newcommand{\inv}{\operatorname{inv}}
\title{A counterexample to signed cycle-restriction positivity\\\large MathDB 350721}
\author{Principia Math In Progress}
\date{September 18, 2026}
\begin{document}
\maketitle
\begin{abstract}
The signed cycle-restriction positivity conjecture fails at $n=8$,
$k=3$, including the corrected formulation in the published paper.
The coefficient of $x_1x_4x_6x_7y_2y_3y_5y_8$ is $-1$.
We prove this by an exhaustive six-row table derived directly from
the required permutation statistics. A second coefficient is $+1$,
so changing a global sign cannot restore positivity.
\end{abstract}

\section*{Task 1. Full solution}
\subsection*{The intended published statement}
For a permutation $\pi\in\mathfrak S_n$, written as
$\pi(1)\cdots\pi(n)$, define
\[
 \RLM(\pi)=\{\pi(j):\pi(j)<\pi(l)\text{ for every }l>j\},
 \qquad
 \EXC(\pi)=\{\pi(j):\pi(j)>j\}.
\]
For a set $S$, put $x_S=\prod_{s\in S}x_s$, and similarly for $y_S$.
Let $c(\pi)$ denote the number of cycles. Write
\begin{equation}\label{eq:P}
 P_{n,k}(\mathbf x,\mathbf y)=(-1)^{n-1}
 \sum_{\substack{\pi\in\mathfrak S_n\\
                  \text{every cycle of }\pi\text{ has length}\ge k}}
 (-1)^{\inv(\pi)}x_{\RLM(\pi)}y_{\EXC(\pi)}.
\end{equation}
Conjecture 5.2.1 of the journal version \cite[p.~411]{final} asserts
$P_{n,k}\in\mathbb N[\mathbf x,\mathbf y]$ for $2\le k\le n$.
The current preprint's Conjecture 40 and the MathDB transcription omit
$(-1)^{n-1}$ and allow $k=1$. The journal corrects both points.
Our counterexample uses the journal's corrected statement.

\begin{theorem}\label{thm:counter}
For $n=8$, $k=3$, let
\[
 M=x_1x_4x_6x_7y_2y_3y_5y_8,\qquad
 N=x_1y_2y_3y_4y_5y_6y_7y_8.
\]
Then $[M]P_{8,3}=-1$ and $[N]P_{8,3}=1$.
In particular, the published conjecture is false, and the
unnormalized polynomial in MathDB also has coefficients of both signs.
\end{theorem}

\subsection*{Deriving the complete table}
To contribute to $[M]P_{8,3}$, a permutation must have
\begin{equation}\label{eq:sets}
 \RLM(\pi)=\{1,4,6,7\},\qquad \EXC(\pi)=\{2,3,5,8\}.
\end{equation}
It must in particular have no fixed points. Let $p_v=\pi^{-1}(v)$.
The excedance conditions for the values $2,3$ force $p_2=1$, $p_3=2$.
The rightmost entry is the largest right-to-left minimum value, so
$\pi(8)=7$. Value $6$ is not an excedance and is not fixed, hence
$p_6>6$. Since position $8$ is occupied, $p_6=7$.
Similarly $p_4>4$. The right-to-left minima occur in increasing
value order, so $p_4<p_6=7$; therefore $p_4\in\{5,6\}$.
Finally $5$ is an excedance value, giving $p_5<5$, and positions
$1,2$ are occupied; thus $p_5\in\{3,4\}$.

We have reduced every possible permutation to
\[
 \pi=2\,3\,\boxed{\quad\quad\quad\quad}\,6\,7,
\]
where the middle four entries are $1,4,5,8$, subject to the conditions
just stated and $p_1<p_4$.
If $p_4=5$, the values $1,5$ fill positions $3,4$ in either order,
and $8$ fills position $6$. If $p_4=6$, place $5$ at position $3$
or $4$, and fill the remaining two positions with $1,8$ in either
order. These give exactly the following six rows. Direct inspection
confirms that each has the two statistic sets in \eqref{eq:sets}.

\begin{center}
\begin{tabular}{llll}
\toprule
One-line form & Disjoint cycles & $c(\pi)$ & Contribution to $[M]P_{8,3}$\\
\midrule
$23154867$ & $(123)(45)(687)$ & $3$ & excluded: a $2$-cycle\\
$23158467$ & $(123)(45876)$ & $2$ & $-1$\\
$23514867$ & $(12354)(687)$ & $2$ & $-1$\\
$23518467$ & $(12358764)$ & $1$ & $+1$\\
$23581467$ & $(1235)(4876)$ & $2$ & $-1$\\
$23851467$ & $(12387645)$ & $1$ & $+1$\\
\bottomrule
\end{tabular}
\end{center}
Indeed $\operatorname{sgn}(\pi)=(-1)^{8-c(\pi)}$, so the weight
including the global factor in \eqref{eq:P} is $(-1)^{c(\pi)-1}$.
All cycles in the last five rows have length at least $3$.
Their contributions add to $-1-1+1-1+1=-1$. This proves the first
coefficient assertion without relying on computer enumeration.

\subsection*{The other sign and the source correction}
For the monomial $N$, all values $2,3,\ldots,8$ must be excedances.
Induction forces $p_v=v-1$ for $2\le v\le8$: the earlier positions
are already occupied by the smaller excedance values. Thus the
unique permutation is $23456781=(12345678)$. Its right-to-left
minimum set is $\{1\}$ and its normalized sign is $+1$.
Therefore $[N]P_{8,3}=1$, completing the proof.\hfill$\square$

For the unnormalized polynomial printed in MathDB, the coefficients
of $M,N$ are $+1,-1$, respectively. Hence this resolves the intended
research assertion even after the journal's corrections; the missing
global sign is not the mathematical obstruction established here.

\subsection*{Verification and attribution}
The conjecture, statistics, and the known $k=2$ result are due to
Alexandersson and Getachew Kebede \cite{final,preprint}. Our six-row
coefficient calculation disproves their proposed extension.
The accompanying exact verifier checks the table and inversion signs,
enumerates all $8!$ permutations, independently generates the allowed
permutations by their cycles, and checks the known derangement
identity and all smaller parameters. These computations supplement
the complete table argument above.

All 23 pages of the current primary, all 27 pages of the journal
final, and all nine pages of the proceedings version were read.
The later refinements by Pei and Zeng \cite{pz} concern the
derangement case; their stated results do not settle the cycle-length
restriction. Bounded public checks found no earlier full claim for
this target. Review is by the authoring agent, with no independent
peer review or proof-assistant certification.

\section*{Tasks 2--5. Proposed follow-up directions}
\begin{enumerate}[label=\textbf{Task \arabic*.},start=2,leftmargin=*]
\item Construct infinite families of negative coefficients for fixed
$k=3$, extending the six-row configuration above.
\item Classify the pairs $(n,k)$ for which all coefficients of
$P_{n,k}$ are nonnegative; the cases $k=2$ and $k>n/2$ are known
or immediate.
\item Determine which identifications of the $x$ and $y$ variables
restore positivity, and whether the fully bivariate specialization
has nonnegative coefficients for every cycle restriction.
\item Describe precisely how forbidding $2$-cycles removes
sign-canceling partners in the source involutions, and derive a
formula for the resulting negative terms.
\end{enumerate}
\begin{thebibliography}{9}
\bibitem{final} P. Alexandersson and F. G. Kebede,
\emph{An involution on derangements preserving excedances and
right-to-left minima}, Australas. J. Combin. \textbf{86}(3) (2023),
387--413. \href{https://ajc.maths.uq.edu.au/pdf/86/ajc_v86_p387.pdf}{Journal PDF}.
\bibitem{preprint} P. Alexandersson and F. G. Kebede, same title,
\href{https://arxiv.org/abs/2105.08455v3}{arXiv:2105.08455v3}, 2022.
\bibitem{pz} Y. Pei and J. Zeng, \emph{Counting derangements with
signed right-to-left minima and excedances},
\href{https://arxiv.org/abs/2206.11236v2}{arXiv:2206.11236v2}, 2023.
\end{thebibliography}
\end{document}
